\scanpage{36}
\mspage{17}
et ceci posé, la relation (56) équivaut à
\[
(60)\qquad\left\{\begin{aligned}
\alpha'&=\xi'\unsure{\rho^{\deleted{-1}}|\rho^{-1}}(\xi')^{-1}\\
\text{i.e.}\quad\eta\sigma_\infty\unsure{(\eta)|\eta)}^{-1}&=(\xi'\unsure{\rho|\rho^{-1}}(\xi')^{-1})\sigma_\infty(\lambda)
\end{aligned}\right.
\]
\deleted{i.e. à}\quad\deleted{\unsure{(61)}}
En effet, la formule (75bis) exprimant $\alpha$ en termes de $\widetilde\alpha$ \unsure{e|$\varepsilon$}
\[
(x)\qquad\alpha=\rho(\eta\widetilde\alpha\rho^{-1}\unsure{(\eta)|\eta)}^{-1})
\]
se récrit, en passant aux inverses et posant
\[
(xx)\qquad\widetilde\alpha'=\widetilde\alpha^{-1}
\]
sous la forme
\[
(61)\qquad\uline{\alpha'}=\rho(\rho^{-1}\unsure{(\eta)|\eta)}\widetilde\alpha'\eta^{-1})
=\uline{\eta\rho(\widetilde\alpha')\rho\unsure{(\eta)|\eta)}^{-1}}
\]
\deleted{et la formule}
(on considère ici $\alpha,\alpha',\widetilde\alpha,\widetilde\alpha'$ comme déduits de $\lambda,\eta\in\widehat\Pi_{0,3}$ par la première relation (55) pour ce qui est de $\widetilde\alpha$, dont les autres se déduisent par $(x),(xx)$. D'ailleurs, la \inserted{2\textsuperscript{e}} formule (55) se\deleted{\unsure{t}} récrit alors, en passant aux inverses
\[
\widetilde\alpha'=\rho^{-1}(\xi)\xi^{-1}
\]
i.e.
\[
\rho(\widetilde\alpha')=\xi\rho(\xi)^{-1}
\]
\deleted{de sorte que (61), \unsure{(62)} s'écrit encore}
\inserted{qui est équivalt \unsure{à (60)}}
\[
\alpha'=\eta\xi\rho(\xi^{-1}\eta^{-1})
\]
\deleted{i.e. \unsure{donc|$\alpha'$} à (60)}\unsure{,|)}
i.e à (60), avec
\[
\xi'=\eta\xi.
\]

\uline{En résumé} Les autom. du groupoïde $\mathrm{VI}_J$, induisant l'identité sur les sommets et \uline{commutant à l'action de $G_J$}, correspondent aux
\[
(62)\qquad\underline{\left|\,(\lambda,\eta,\xi')\in\widehat\Pi_{03}\times\widehat\Pi_{03}\times\widehat\Pi_{03}\,\right\}}
\]
satisfaisant la relation des lacets \unsure{?|:}
\margin{35}
